% T12 C/D, chapitre 2 : Probabilités. Source : le livre « Terminales scientifiques » (Proba/Lesona, Proba/Fiche, Proba/Exo).
% Les exercices tirés d’autres ouvrages ne sont pas repris (décision du 2 novembre 2026).

\TCDLessonStyleOn
\chapter{Probabilit\'e}
\newcommand{\card}[1]{\text{Card }#1}
\newcommand{\A}[2]{A_{#1}^{#2}}
\newcommand{\C}[2]{C_{#1}^{#2}}
      \section{D\'enombrement}
            %
            \begin{defi}[Cardinal d'un ensemble fini]
\index{Cardinal}
            On appelle \emph{cardinal} d'un ensemble fini le nombre de ses \'el\'ements.
            Si $X$ est un ensemble alors nous noterons par $Card(X)$ le cardinal de $X$.
            Si on se place dans un ensemble $E$ et si $A$ est un sous-ensemble de $E$, alors on note par $\bar{A}$ le compl\'ementaire de $A$ dans $E$, c'est-\`a-dire l'ensemble des \'el\'ements de $E$ qui ne sont pas dans $A$. 
            \end{defi}
            %
            \begin{prp}
            Soient $A$ et $B$ deux ensembles finis. On a alors :
            \begin{itemize}
                  \item Si $A$ et $B$ sont deux ensembles disjoints, alors $Card(A \cup B) = Card(A) + Card(B)$.
                  \item $Card(A \cup B) = Card(A) + Card(B) - Card(A \cap B)$.
                  \item Si $\bar{A}$ d\'esigne le compl\'ementaire de $A$ dans $E$, alors $Card(\bar{A}) = Card(E) - Card(A)$.
            \end{itemize}
            \end{prp}
            %
            Soient $E$ et $F$ deux ensembles finis non-vides \`a $p$ et $n$ \'el\'ements respectivement.
            Notons par $e_1,\ e_2,\ \ldots\ ,\ e_p$ les \'el\'ements de $E$ 
            et par  $f_1,\ f_2,\ \ldots\ ,\ f_n$ les \'el\'ements de $F$. 
            %
            \subsection{Nombre d'applications de $E$ vers $F$}
            Si $\varphi : E \rightarrow F $ 
            est une application, alors :
            \begin{itemize}
                  \item il y a $n$ possibilit\'es pour $\varphi(e_1)$,
                  \item il y a $n$ possibilit\'es pour $\varphi(e_2)$,
                  \item et ainsi de suite,
                  \item il y a $n$ possibilit\'es pour $\varphi(e_p)$.
            \end{itemize}
            Ainsi, il y a exactement \propbox{$n^p$} applications possibles de $E$ dans $F$.
            %
            \begin{figure}[h!]
            \centering
                  \begin{tikzpicture}[line cap=round,line join=round,>=triangle 45,x=1.0cm,y=1.0cm]
                  \clip(-2.7,0.26) rectangle (5.44,5.32);
                  \draw [rotate around={89.01:(-1.56,2.7)}] (-1.56,2.7) ellipse (1.38cm and 0.74cm);
                  \draw [rotate around={89.16:(4.25,2.72)}] (4.25,2.72) ellipse (2.27cm and 1cm);
                  \draw [->] (-1.82,3.63) -- (4.64,3.75);
                  \draw [->] (-1.37,1.57) -- (4.64,3.75);
                  \draw [->] (-1.42,3) -- (4.43,2.6);
                  \draw [->] (-1.81,2.24) -- (4.44,1.05);
                  \begin{scriptsize}
                  \fill [color=black] (-1.82,3.63) circle (1.5pt);
                  \draw[color=black] (-1.67,3.77) node {$e_1$};
                  \fill [color=black] (-1.42,3) circle (1.5pt);
                  \draw[color=black] (-1.27,3.13) node {$e_2$};
                  \fill [color=black] (-1.81,2.24) circle (1.5pt);
                  \draw[color=black] (-1.64,2.37) node {$e_3$};
                  \fill [color=black] (-1.37,1.57) circle (1.5pt);
                  \draw[color=black] (-1.25,1.75) node {$e_4$};
                  \fill [color=black] (3.98,4.42) circle (1.5pt);
                  \draw[color=black] (4.11,4.56) node {$f_1$};
                  \fill [color=black] (4.64,3.75) circle (1.5pt);
                  \draw[color=black] (4.79,3.88) node {$f_2$};
                  \fill [color=black] (3.66,3.09) circle (1.5pt);
                  \draw[color=black] (3.8,3.23) node {$f_3$};
                  \fill [color=black] (4.43,2.6) circle (1.5pt);
                  \draw[color=black] (4.46,2.79) node {$f_4$};
                  \fill [color=black] (3.89,1.71) circle (1.5pt);
                  \draw[color=black] (4.03,1.85) node {$f_5$};
                  \fill [color=black] (4.44,1.05) circle (1.5pt);
                  \draw[color=black] (4.58,1.19) node {$f_6$};
                  \end{scriptsize}
                  \end{tikzpicture}
                  \caption{
                           Exemple d'application d'un ensemble $E$ de $4$ \'el\'ements 
                           vers un ensemble $F$ de $6$ \'el\'ements.
                           }
            \end{figure}
            \subsection{Arrangements dans un ensemble fini}
\index{Arrangement}
                  \begin{defi}
                        Soient $E$ un ensemble fini \`a $n$ \'el\'ements et $p$ un nombre entier inf\'erieur ou \'egal \`a $n$.
                        On appelle \emph{arrangement de $p$ \'el\'ements de $E$}, et on note 
                        \defbox{$\displaystyle{\Arr{n}{p}}$}, le nombre de mani\`eres d'ordonner $p$ \'el\'ements de $E$.
                  \end{defi}
                  %
                  	\begin{prop}
                  			\propbox{
                 				\[
                  					\Arr{n}{p} = \frac{n!}{(n-p)!}
                  				\]}
                  	\end{prop}
						%
						\begin{ex}
							Consid\'erons quatre personnes num\'erot\'ees de 1 \`a 4
							ainsi que six chaises marqu\'ees par les lettres a, b, \ldots, f.
							%
							On veut conna\^itre le nombre de fa\c{c}ons dont ces quatre personnes
							peuvent s'asseoir sur les six chaises ;
							on ne tol\`ere pas que deux personnes s'assoient sur la m\^eme chaise.
							\begin{valiny}
								Pour cel\`a, posons	$E = \{ 1 ,\ 2,\ 3,\ 4 \}$ et 
															$F = \{ a,\ b,\ c,\ d,\ e,\ f \}$.
								\begin{itemize}
									\item La personne num\'ero $1$ a $6$ choix possibles et, apr\`es qu'elle a pris sa place,
									\item la personne num\'ero $2$ a $5$ choix possibles,
									\item la personne num\'ero $3$ a $4$ choix possibles,
									\item la personne num\'ero $4$ a $3$ choix possibles.
								\end{itemize}
								En tout, cel\`a fait $6\times 5 \times 4 \times \cdots \times 3 = 360$ possibilit\'es.
								%
									Ce nombre peut aussi se calculer par la formule
								$\A{6}{2} 
									= \frac{6!}{2!} 
									= \frac{6\times 5 \times 4 \times \cdots \times 3 \times 2 \times 1}{2 \times 1}
									= 360
								$.
							\end{valiny}
						\end{ex}
					%
            \subsection{Permutation dans un ensemble fini}
\index{Permutation}
                  Si on prend $E = F$ dans le cas pr\'ec\'edent, alors on se retrouve avec ce que l'on appelle
                  \emph{permutation de $E$}.
                  Le nombre de permutations d'un ensemble \`a $n$ \'el\'ements est $n!$.
                  \begin{ex}
                        Consid\'erons une course avec six participants. 
                        A l'issue de la course, on suppose qu'il n'y a pas d'ex aequo.
                        On veut d\'eterminer le nombre de possibilit\'es de r\'esultats pour une telle course.
                     \begin{valiny}
                        \begin{itemize}
                              \item Il y a $6$ possibilit\'es pour la premi\`ere place,
                              \item il y a $5$ possibilit\'es pour la seconde place,
                              \item $\vdots$ 
                              \item il y a $1$ possibilit\'e pour la derni\`ere place.
                        \end{itemize}
                        Ainsi, il y a $6 \times 5 \times \cdots \times 1$ issues possibles.
                        Ce nombre est $6! = 720$.
                     \end{valiny}
                  \end{ex}
				\subsection{Combinaison dans un ensemble fini}
\index{Combinaison}
					\begin{defi}
						Soient $E$ un ensemble fini \`a $n$ \'el\'ements et $p$ un nombre entier naturel inf\'erieur ou \'egal \`a $n$.
						On appelle \emph{nombre de combinaisons de $p$ \'el\'ements de $E$} le nombre, not\'e 
						\defbox{$\displaystyle{\Comb{n}{p}}$},
							de mani\`eres de choisir $p$ \'el\'ements de $E$.
					\end{defi}
					%
                  	\begin{prop}~\\[-1\baselineskip]
								\propbox{
									\[
										\Comb{n}{p} = \frac{n!}{p!\;(n-p)!}
									\]
								}
 						\end{prop}
 					%
 					\begin{ex}
	 						On dispose d'un d\'e cubique \`a six faces, num\'erot\'ees de $1$ \`a $6$.
 						On veut colorer deux des faces de ce d\'e, de mani\`ere \`a cacher les num\'eros y figurant.
 						De combien de mani\`eres peut-on le faire\,?
 						\begin{valiny}
	 						Il suffit de << choisir >> deux faces parmi les 6.
 							Il y a, par cons\'equent, 
 							$\Comb{6}{2} = \frac{6!}{2! \; (6-2)!} = 15$ fa\c{c}ons de colorer deux faces d'un d\'e cubique.
 						\end{valiny}
 					\end{ex}
%
      \section{Notion de probabilit\'e}
            \begin{defi}[\'Epreuve ou exp\'erience al\'eatoire]
\index{Epreuve@\'Epreuve}
\index{Experience aleatoire@Exp\'erience al\'eatoire}
                  Une \emph{exp\'erience al\'eatoire} est une exp\'erience renouvelable
                  dont le r\'esultat ne peut \^etre pr\'evu, et qui, renouvel\'ee dans des conditions identiques
                  ne donne pas forc\'ement le m\^eme r\'esultat \`a chaque renouvellement.
                  \begin{ex}
                        \begin{itemize}
                              \item Lancement d'un d\'e,
                              \item tirage de boules dans une urne.
                        \end{itemize}
                  \end{ex}
            \end{defi}
            %
            \begin{defi}[Univers, \'eventualit\'es, \'ev\`enement]
\index{Univers des eventualite@Univers des \'eventualit\'es}
                  L'ensemble de tous les r\'esultats possibles est appel\'e \emph{univers des \'eventualit\'es}.
                  L'univers est souvent not\'e $\Omega$.
                  %
                  Les \'el\'ements de $\Omega$ sont appel\'es \emph{\'eventualit\'es}.
\index{Eventualite@\'Eventualit\'e}
                  On appelle \emph{\'ev\`enement} tout ensemble d'\'eventualit\'es, c'est-\`a-dire toute partie de $\Omega$.
\index{Evenements@\'Ev\`enements}
                  Si $A$ est un \'ev\`enement, on note par $\bar{A}$ le compl\'ementaire de $A$ dans $\Omega$,
                  $\bar{A}$ est appel\'e l'\emph{\'ev\`enement contraire de $A$}.
\index{Evenements@\'Ev\`enements!contraires}
                  Deux \'ev\`enements $A$ et $B$ sont dits \emph{incompatibles} si $A\cap B = \emptyset$.
\index{Evenements@\'Ev\`enements!incompatibles}
                  \begin{ex}
                        Consid\'erons les \'epreuves suivantes :
                        \begin{enumerate}
                              \item Lancer d'un d\'e \`a six faces.
                                    L'univers des possibilit\'es est 
                                    $\Omega = \{ 1,\ 2,\ 3,\ 4,\ 5,\ 6 \}$.
                              \item Lancer de deux pi\`eces de monnaie distinctes.
                                    L'univers des possibilit\'es est
                                    \[
                                      \Omega = \left\{
                                        \begin{array}{cc}
                                          (\text{Pile},\text{Pile}), & (\text{Pile},\text{Face}), \\
                                          (\text{Face},\text{Pile}), & (\text{Face},\text{Face})
                                        \end{array}
                                      \right\}.
                                    \]
                        \end{enumerate}
                  \end{ex}
                  %
                  \begin{ex}
                        On lance un d\'e cubique dont les six faces sont num\'erot\'ees par les nombres 
                        1, 2, 3, 4, 5 et 6.
                        
                        \begin{enumerate}
                              \item L'univers est : $\Omega = \{1,\ 2,\ 3,\ 4,\ 5,\ 6 \}$.
                              \item Consid\'erons les \'ev\`enements : 
                                    \begin{align*}
                                    A &: \text{<< Obtenir un nombre pair >>} \\
                                    B &: \text{<< Obtenir un nombre inf\'erieur ou \'egal \`a 3 >>} .
                                    \end{align*}
                                    On a donc $A = \{ 2,\ 4,\ 6 \}$ et $B = \{ 1,\ 2,\ 3 \}$.
                              \item L'\'ev\`enement << $A$ ou $B$ >> est $A \cup B = \{ 1,\ 2,\ 3,\ 4,\ 6 \}$.
                              \item L'\'ev\`enement << $A$ et $B$ >> est $A \cap B = \{ 2 \}$, 
                                    c'est un \emph{\'ev\`enement \'el\'ementaire}.
                              \item $\bar{A} = \{ 1,\ 3,\ 5 \}$ et $\bar{B} = \{ 4,\ 5,\ 6 \} $.
                              \item $A$ et $\bar{A}$ sont incompatibles.
                                    En fait, quel que soit l'\'ev\`enement $X$, $X$ et $\bar{X}$ sont toujours incompatibles.
\index{Evenements@\'Ev\`enements!elementaires@\'el\'ementaires}
                        \end{enumerate}
                  \end{ex}
            \end{defi}
      Consid\'erons un univers fini $\Omega = \{ \omega_1,\ \omega_2,\ \ldots ,\omega_n \}$.
      \begin{defi}[Probabilit\'e]
\index{Probabilite@Probabilit\'e}
            On appelle \emph{probabilit\'e} sur $\Omega$ toute application 
            de l'ensemble des \'ev\`enements dans $[0 ; 1]$ telle que : 
            \begin{itemize}
                  \item la somme des probabilit\'es des \'ev\`enements \'el\'ementaires est \'egale \`a $1$,
                  \item la probabilit\'e d'un \'ev\`enement est \'egale \`a la somme des probabilit\'es
                        des \'ev\`enements \'el\'ementaires qui le constituent.
            \end{itemize}
      \end{defi}
      %
      \begin{prop}
            Soient $\Omega$ l'univers et $A$ et $B$ deux \'ev\`enements.
            Soit $P$ une probabilit\'e sur $\Omega$. On a :
            \begin{itemize}
                  \item $P(\emptyset ) = 0$ et $P(\Omega ) = 1$,
                  \item $P(\bar{A}) = 1 - P(A)$,
                  \item $P(A\cup B) = P(A) + P(B) - P(A\cap B)$.
                        Si $A$ et $B$ sont incompatibles, alors $P(A\cup B) = P(A) + P(B)$.
                  \item Si $A\subseteq B$ alors $P(A) \leq P(B)$.
            \end{itemize}
      \end{prop}
      %
      \subsection*{Le cas \'equiprobable}
\index{Equiprobabile@\'Equiprobable}
            On dit qu'il y a \emph{\'equiprobabilit\'e} lorsque tous les \'ev\`enements 
            \'el\'ementaires ont la m\^eme probabilit\'e.
            Dans le cas \'equiprobable, si $\omega$ est un \'ev\`enement \'el\'ementaire, alors
                  $P(\omega) = \frac{1}{\card{\Omega}}$.
            Si $A$ est un \'ev\`enement quelconque, alors 
            \thmbox{
                  $$P(A) = \frac{\card{A}}{\card{\Omega}} 
                        = \frac{\text{Nombre de cas favorables}}{\text{Nombre de cas possibles}}.$$}
            \begin{exe}
                        On consid\`ere un lancer de d\'e non truqu\'e.
                        La probabilit\'e de chaque \'ev\`enement \'el\'ementaire est \'egale \`a $\frac{1}{6}$.
                        Soient les \'ev\`enements : 
                              \begin{enumerate}
                                    \item A : << Obtenir un nombre pair. >>
                                    \item B : << Obtenir un nombre inf\'erieur ou \'egal \`a 3. >>
                              \end{enumerate}
                        Calculer les probabilit\'es : $P(A)$, $P(B)$, $P(A\cap B)$ et $P(A \cup B)$.
            \end{exe}
            %
            \begin{exe}
                        On consid\`ere maintenant un d\'e truqu\'e de sorte que 
                              la probabilit\'e d'obtenir $6$ soit deux fois la probabilit\'e d'obtenir $1$,
                              la probabilit\'e d'obtenir $1$ \'etant la m\^eme que celle d'obtenir $2$, $3$, $4$ ou $5$.
                              \begin{enumerate}
                                    \item Calculer la probabilit\'e de chaque \'ev\`enement \'el\'ementaire.
                                    \item Calculer la probabilit\'e des \'ev\`enements :
                                          \begin{enumerate}
                                                \item A : << Obtenir un nombre pair. >>
                                                \item B : << Obtenir un nombre inf\'erieur ou \'egal \`a 3. >>
                                          \end{enumerate}
                              \end{enumerate}
            \end{exe}
            %
            \begin{exe}
                        Dans une urne contenant 10 boules num\'erot\'ees de 0 \`a 9,
                        on tire successivement et avec remise 3 boules.
                        \begin{enumerate}
                              \item Calculer le nombre de tirages possible.
                              \item Calculer le nombre de tirages telle que 
                                    la premi\`ere boule ne soit pas celle num\'erot\'ee 0.
                              \item Calculer la probabilit\'e que les num\'eros tir\'es dans l'ordre
                                    forment un nombre de trois chiffres
                                    (c'est-\`a-dire ne commen\c{c}ant pas par 0).
                        \end{enumerate}
                        (On admet que toutes les boules ont la m\^eme probabilit\'e d'\^etre tir\'ee.)
            \end{exe}
            %
            \begin{exe}
                        On consid\`ere un urne contenant 26 boules,
                        chacune marqu\'ee d'une lettre de l'alphabet (fran\c{c}aise) diff\'erente.
                        On tire successivement et sans remise cinq boules de l'urne.
                        \begin{enumerate}
                              \item Calculer le nombre de tirages possible.
                              \item Calculer la probabilit\'e d'obtenir les cinq lettres `M A T H S' dans cet ordre.
                              \item Calculer la probabilit\'e d'obtenir les cinq lettres `A H M S T',
                                    pas n\'ecessairement dans cet ordre.
                              \item Calculer la probabilit\'e d'obtenir cinq voyelles.
                        \end{enumerate}
            \end{exe}
            %
            \begin{exe}
                        Dans un jeu de 52 cartes, on tire 13 cartes simultan\'ement.
                        Calculer la probabilit\'e d'avoir : 
                        \begin{enumerate}
                              \item Exactement un as,
                              \item exactement deux as,
                              \item exactement trois as,
                              \item exactement quatre as,
                              \item exactement cinq as,
                              \item uniquement des tr\`efles.
                        \end{enumerate}
            \end{exe}
      %
      %
      \section{Probabilit\'e conditionnelle}
\index{Probabilite@Probabilit\'e!conditionnelle}
            \subsection{Probabilit\'e conditionnelle}
            Consid\'erons un univers $\Omega$ et soient $A$ et $B$ deux \'ev\`enements.
            Soit $P$ une probabilit\'e.
            On suppose que $P(B) \neq 0$.
            La \emph{probabilit\'e de $A$ sachant $B$} est 
                  \defbox{
                              \[
                                    P(A|B) = \frac{P(A\cap B)}{P(B)}.
                              \]
                         }
            \begin{ex}
                  Une urne contient neuf boules indiscernables au toucher.
                  Les boules sont num\'erot\'ees de 1 \`a 9.
                  On tire une boule au hasard.
                  Soient les \'ev\`enements : 
                  \begin{enumerate}
                        \item A : << Le num\'ero tir\'e est pair. >>
                        \item B : << Le num\'ero tir\'e est un multiple de 3. >>
                  \end{enumerate}
                  
                  On a $ \Omega = \{ 1,\ 2,\ 3,\ 4,\ 5,\ 6,\ 7,\ 8,\ 9 \}$,
                       $ A = \{ 2,\ 4,\ 6,\ 8 \}$ et
                       $ B = \{ 3,\ 6,\ 9 \}$.
                  Il suit que $A \cap B = \{ 6 \}$.
                  Ainsi, $P(A) = \frac{4}{9}$, $P(B) = \frac{3}{9} = \frac{1}{3}$ et $P(A\cap B) = \frac{1}{9}$.
                  
                  Supposons qu'on sait qu'on a obtenu un multiple de 3.
                  Quelle est alors la probabilit\'e pour que ce soit un num\'ero pair\,?
                  
                  Il suffit de calculer 
                        \[
                        P(A|B) = \frac{P(A\cap B)}{P(B)} = \frac{\frac{1}{9}}{\frac{1}{3}} = \frac{1}{3}.
                        \]
            \end{ex}
            \smallnota{
            La probabilit\'e de $A$ sachant $B$ est aussi parfois not\'e $P_B (A)$.
            }
            %
            \begin{exe}
                  On vous propose de jouer \`a un jeu.
                  On pose devant vous cinq paniers clos ; 
                  l'un d'eux contient une pi\`ece de monnaie tandis que
                  les quatre autre sont vides.
                  %
                  Votre objectif est de trouver la pi\`ece de monnaie.
                  
                  Vous avez le droit de choisir un panier, 
                  vous n'avez aucun moyen de d\'ecider d'une mani\`ere rationnelle quel panier choisir.
                  %
                  Apr\`es que vous avez fait votre choix, on ouvre trois des paniers qui s'av\`erent \^etre vides ;
                  votre panier reste clos.
                  %
                  Il reste donc deux paniers clos : le v\^otre et un autre.
                  
                  On vous demande ensuite votre d\'ecision finale : 
                        allez-vous changer votre choix et miser sur l'autre panier ou
                        allez-vous garder votre choix initial\,?
                        
                        \paragraph{Illustration}
                              Supposons que les cinq paniers sont num\'erot\'es de 1 \`a 5.
                              
                                    \begin{tikzpicture}[scale=2]
                                    \foreach \x in {1,2,3,4,5}
                                          \draw[pattern=checkerboard,pattern color=loko2] 
                                                (\x-.3,.3) -- (\x+.3,.3) 
                                                {[rounded corners = 10] -- (\x+.3,-.3) -- (\x-.3,-.3) -- (\x-.3,.3)}
                                                (\x,0) node[fill=white]{\x}
 ;
                                    \end{tikzpicture}
                                    
                              Disons que vous avez choisi le panier 1.
                              
                                    \begin{tikzpicture}[scale=2]
                                    \foreach \x in {1,2,3,4,5}
                                          \draw[pattern=checkerboard,pattern color=loko2] 
                                                (\x-.3,.3) -- (\x+.3,.3) 
                                                {[rounded corners = 10] -- (\x+.3,-.3) -- (\x-.3,-.3) -- (\x-.3,.3)}
                                                (\x,0) node[fill=white]{\x}
 ;
                                          \draw[color=loko1] (1,0) circle (.5);
                                    \end{tikzpicture}
                                    
                              On ouvre ensuite les paniers 2, 3 et 4 qui s'av\`erent \^etre vides.
                              
                                    \begin{tikzpicture}[scale=2]
                                    \foreach \x in {1,2,3,4,5}
                                          \draw[pattern=checkerboard,pattern color=loko2] 
                                                (\x-.3,.3) -- (\x+.3,.3) 
                                                {[rounded corners = 10] -- (\x+.3,-.3) -- (\x-.3,-.3) -- (\x-.3,.3)}
                                                (\x,0) node[fill=white]{\x}
 ;
                                          \draw[color=loko1] (1,0) circle (.5);
                                    \foreach \x in {2,3,4}
                                          \draw[line width=5,color=loko1,opacity=.5] 
                                                (\x+.3,.3) -- (\x-.3,-.3) 
                                                (\x+.3,-.3) -- (\x-.3,.3)
 ;
                                    \end{tikzpicture}
                                    
                              Il reste alors deux paniers clos : les num\'eros 1 et 5.
                              
                                    \begin{tikzpicture}[scale=2]
                                    \foreach \x in {1,5}
                                          \draw[pattern=checkerboard,pattern color=loko2] 
                                                (\x-.3,.3) -- (\x+.3,.3) 
                                                {[rounded corners = 10] -- (\x+.3,-.3) -- (\x-.3,-.3) -- (\x-.3,.3)}
                                                (\x,0) node[fill=white]{\x}
 ;
                                    \foreach \x in {1,5}
                                          \draw[opacity=.5] 
                                                (\x,0) node[fill=white]{{\fontsize{40}{40} \selectfont\,?} }
 ;
                                    \end{tikzpicture}
                                    
                              Allez-vous changer votre choix et miser sur le num\'ero 5
                              ou allez-vous rester ferme \`a votre intuition initiale et garder le num\'ero 1\,?
            \end{exe}
            %
            \subsection{\'Ev\`enements ind\'ependants}
\index{Evenements@\'Ev\`enements!independants@ind\'ependants}
            Deux \'ev\`enements $A$ et $B$ sont dits \emph{ind\'ependants} si 
                        \defbox{$$P(A\cap B) = P(A) \cdot P(B).$$}
            Si deux \'ev\`enements $A$ et $B$ sont ind\'ependants, on a :
                  \begin{displaymath}
                        P(A|B) = \frac{P(A\cap B)}{P(B)} = \frac{P(A) \cdot P(B)}{P(B)} = P(A).
                  \end{displaymath}
            %
            \begin{ex}
                  On dispose d'une pi\`ece de monnaie non truqu\'ee.
                  On proc\`ede \`a deux lancers successifs.
                  Calculer la probabilit\'e d'obtenir pile puis face.
                  \begin{valiny}
                        Soient les deux \'ev\`enements : 
                              \begin{itemize}
                                    \item A : << Le premier lancer donne pile. >>
                                    \item B : << Le second lancer donne face. >>
                              \end{itemize}
                        Ces deux \'ev\`enements sont ind\'ependants --- 
                        le r\'esultat de la premi\`ere n'affecte pas celui de la seconde.
                        On cherche $P(A\cap B)$.
                        On sait que 
                        $P(A\cap B) = P(A) \cdot P(B) = \frac{1}{2} \cdot \frac{1}{2} = \frac{1}{4}$.
                  \end{valiny}
            \end{ex}
            %
	\section{\'Epreuve de Bernoulli et distribution binomiale}
                  \subsection{\'Epreuve de Bernoulli}
\index{Bernoulli}
                        Dans certaines exp\'eriences al\'eatoires, 
                        on s'interesse \`a la r\'ealisation d'un \'ev\`enement.
                        On parlera de << \emph{succ\`es} >> si l'\'ev\`enement se r\'ealise
\index{succes@succ\`es}
                        et d'<<~\emph{\'echec} >> sinon.
\index{echec@\'echec}
                        Ce type d'exp\'erience est appel\'e \emph{\'epreuve de Bernoulli}.
\index{Bernoulli!epreuve de@\'epreuve de}
                        On appelle << \emph{sch\'ema de Bernoulli} >> une succession d'\'epreuves ind\'ependantes de ce type.
\index{Bernoulli!schema de@sch\'ema de}
                        %
                        \begin{ex}
                              On lance $n$ fois une pi\`ece,
                              on consid\`ere comme un succ\`es le fait d'obtenir pile
                              (et donc comme \'echec si on obtient face).
                              
                              Le tirage successif et avec remise d'une boule dans une urne.
                              Le succ\`es peut \^etre d\'efini comme l'obtention d'un certain type de boule.
                        \end{ex}
                  \subsection{Distribution binomiale}
                        Consid\'erons un sch\'ema de Bernoulli qui 
                        consiste \`a la succession de $n$ \'epreuves identiques et ind\'ependantes.
                        Soient $p$ la probabilit\'e d'un succ\`es et $q ( = 1-p)$ la probabilit\'e d'un \'echec.
                        Soit $k$ un nombre entier naturel compris entre $0$ et $n$.
                        On veut conna\^itre la probabilit\'e d'avoir $k$ succ\`es au cours des $n$ lancers.
                        
                        La probabilit\'e d'avoir 
                        uniquement des succ\`es au cours des $k$ premiers lancers
                        et uniquement des \'echecs au cours des $n-k$ autres lancers est : 
                              \begin{displaymath}
                                    p^k \cdot (1-p)^{n-k}.
                              \end{displaymath}
                        Il y a $\C{n}{k}$ mani\`eres de r\'epartir les $k$ succ\`es parmi les $n$ \'epreuves.
                        La probabilit\'e d'obtenir $k$ succ\`es est donc 
                                    \thmbox{$$ \C{n}{k}\ p^k \cdot (1-p)^{n-k}. $$}
                        %
                        \begin{ex}
                              On jette 100 fois une pi\`ece de monnaie truqu\'ee de mani\`ere \`a ce que 
                              la probabilit\'e d'obtenir face soit trois fois celle d'obtenir pile. \newline
                              Calculer la probabilit\'e d'avoir 25 fois pile.
                              %
                              \begin{valiny}
                                    Soient $p$ la probabilit\'e d'obtenir pile et $q = 1-p$ 
                                    la probabilit\'e d'obtenir face.
                                    On a $q = 3 p $ c'est-\`a-dire $1-p = 3p$.
                                    En r\'esolvant, on a $p = \frac{1}{4}$ et $q = \frac{3}{4}$.
                                    Les r\'esultats de chaque lancer sont ind\'ependants.
                                    \footnote{Quand l'\'enonc\'e ne pr\'ecise pas,
                                    il est sous-entendu que les lancers sont ind\'ependants.}
                                    On est donc en pr\'esence d'un sch\'ema de Bernoulli.
                                    La probabilit\'e cherch\'ee est :
                                    \begin{displaymath}
                                          P(k = 25) 
                                                = \C{100}{25}\ \left(\frac{1}{4}\right)^{25} \cdot \left(\frac{3}{4}\right)^{75}
                                                = \C{100}{25} \frac{3^{75}}{4^{100}}.
                                    \end{displaymath}
                              \end{valiny}
                        \end{ex}
      \section{Variable al\'eatoire}
            \subsection{Notion de variable al\'eatoire}
                  \begin{ex}
                        Consid\'erons l'\'epreuve qui consiste \`a lancer une pi\`ece trois fois successivement.
                        On note par $p$ la probabilit\'e d'obtenir pile.
                        Soit $X$ le nombre de piles obtenues lors des trois lancers.
                        On appelle $X$ une variable al\'eatoire et on a :
                              \begin{align*}
                                    P(X = 0) & = \C{3}{0}\ p^0 \cdot (1-p)^3 \\
                                    P(X = 1) & = \C{3}{1}\ p^1 \cdot (1-p)^2 \\
                                    P(X = 2) & = \C{3}{2}\ p^2 \cdot (1-p)^1 \\
                                    P(X = 3) & = \C{3}{3}\ p^3 \cdot (1-p)^0 .
                              \end{align*}
                  \end{ex}
                  
                  Plus g\'en\'eralement, si $\Omega$ est l'univers, 
                  une variable al\'eatoire est une application qui \`a chaque \'eventualit\'e associe un nombre r\'eel.
                  
                  L'\emph{univers image} de la variable al\'eatoire $X$ est 
\index{Univers image}
                  l'ensemble des images des \'el\'ements de $\Omega$ par $X$.
                  
                  Pour tout \'el\'ement $k$ de l'univers image, l'ensemble des ant\'ec\'edents de $k$
                  (c'est-\`a-dire l'ensemble des $\omega \in \Omega$ tels que $X(\omega) = k$)
                  par $X$ forme un \'ev\`enement qu'on note $X=k$.
                  
                  La loi de probabilit\'e de $X$ est la donn\'ee des probabilit\'es de tous les \'ev\`enements $(X=k)$.
\index{Loi de probabilite@Loi de probabilit\'e}
                  %
                  \subsection{Fonction de r\'epartition}
\index{Fonction de repartition@Fonction de r\'epartition}
                        Soit $X$ une variable al\'eatoire sur un univers $\Omega$.
                        La \emph{fonction de r\'epartition} $F$ de $X$ est la fonction qui 
                        \`a tout r\'eel $x$ associe $F(x) = P(X \leq x)$.
                        \begin{ex}
                              Consid\'erons l'\'epreuve qui consiste \`a lancer trois fois successivement une pi\`ece normale.
                              Soit $X$ la variable al\'eatoire qui \`a chaque \'eventualit\'e associe le 
                              nombre de piles obtenues.
                              On a alors : 
                              \[
                                    \begin{array}{|c|c|c|c|c|}
                                          \hline
                                          k       &       0       &       1       &       2       &       3       \\
                                          \hline
                                          P(X=k)  &  \frac{1}{8}  &  \frac{3}{8}  &   \frac{3}{8} &  \frac{1}{8}  \\
                                           \hline
                                    \end{array}
                              \]
                              ce qui nous donne
                              \[
                                    \begin{array}{|c|c|c|c|c|c|}
                                          \hline
                                          x    & ]-\infty ; 0[ & [0;1[       & [1;2[       & [2;3[       & [3;+\infty[       \\
                                          \hline
                                          F(x) &       0       & \frac{1}{8} & \frac{1}{2} & \frac{7}{8} &       1           \\
                                          \hline
                                    \end{array}
                              \]
                              \begin{figure}[H]
                              \centering
                                    \begin{tikzpicture}[line cap=round,line join=round,>=triangle 45,x=1.0cm,y=5.0cm]
                                    \draw [color=black, xstep=1.0cm,ystep=1.0cm, line width=.4, opacity=.4] 
                                          (-2.41,-0.1) grid (6.52,1.12);
                                    \draw [color=black, xstep=0.5cm,ystep=0.5cm, line width=.2, opacity=.2] 
                                          (-2.41,-0.1) grid (6.52,1.12);
                                    \draw [color=black, xstep=0.1cm,ystep=0.1cm, line width=.15, opacity=.15] 
                                          (-2.41,-0.1) grid (6.52,1.12);
                                    \draw[->,color=black] (-2.41,0) -- (6.52,0);
                                    \foreach \x in {-2,-1,1,2,3,4,5,6}
                                    \draw[shift={(\x,0)},color=black] (0pt,2pt) -- (0pt,-2pt) node[below] {\footnotesize $\x$};
                                    \draw[->,color=black] (0,-0.1) -- (0,1.12);
                                    \foreach \y in {,0.2,0.4,0.6,0.8,1}
                                    \draw[shift={(0,\y)},color=black] (2pt,0pt) -- (-2pt,0pt) node[left] {\footnotesize $\y$};
                                    \draw[color=black] (0pt,-10pt) node[right] {\footnotesize $0$};
                                    \clip(-2.41,-0.1) rectangle (6.52,1.12);
                                    \draw [line width=2pt,domain=-2.4088237372410504:0.0] plot(\x,{(-0-0*\x)/-2});
                                    \draw [line width=2pt] (0,0.13)-- (1,0.13);
                                    \draw [line width=2pt] (1,0.5)-- (2,0.5);
                                    \draw [line width=2pt] (2,0.88)-- (3,0.88);
                                    \draw [line width=2pt,domain=3.0:6.518929819346457] plot(\x,{(--1-0*\x)/1});
                                    \begin{scriptsize}
                                    \draw [color=black, fill=white] (0,0) circle (1.5pt);
                                    \fill [color=black] (0,0.13) circle (1.5pt);
                                    \draw [color=black, fill=white] (1,0.13) circle (1.5pt);
                                    \fill [color=black] (1,0.5) circle (1.5pt);
                                    \draw [color=black, fill=white] (2,0.5) circle (1.5pt);
                                    \fill [color=black] (2,0.88) circle (1.5pt);
                                    \draw [color=black, fill=white] (3,0.88) circle (1.5pt);
                                    \fill [color=black] (3,1) circle (1.5pt);
                                    \end{scriptsize}
                                    \end{tikzpicture}
                                 \caption{Repr\'esentation graphique de la fonction de r\'epartition.}
                              \end{figure}
                              On constate que :
                              \begin{itemize}
                                    \item $F$ est croissante,
                                    \item $F$ est une fonction en escalier,
                                    \item $\lim_{x\rightarrow-\infty} F(x) = 0$ et $\lim_{x\rightarrow+\infty} F(x) = 1$.
                              \end{itemize}
                        \end{ex}
                  %
                  \subsection{Caract\'eristiques d'une variable al\'eatoire}
                        \subsubsection{\'Esp\'erance math\'ematique}
\index{Esperance mathematique@\'Esp\'erance math\'ematique}
                              Soit $X$ une variable al\'eatoire prenant les valeurs 
                              $x_1,\ x_2,\ \ldots ,\ x_n$ avec les probabilit\'es 
                              $p_1,\ p_2,\ \ldots ,\ p_n$.
                              L'\emph{\'esp\'erance math\'ematique} de $X$ est
                              \par\smallskip\noindent
                                    \defbox{
                                          $\displaystyle
                                          \begin{aligned}
                                                E[X] &= \sum_{i=1}^{n} x_i\cdot p_i \\
                                                     &= x_1\cdot p_1 + x_2\cdot p_2 + \cdots + x_n\cdot p_n.
                                          \end{aligned}$
                                    }
                              \par\smallskip
                              %
                              \begin{ex}
                                    Dans l'exemple pr\'ec\'edent :
                                    $E[X] =  0 \cdot \frac{1}{8}
                                           + 1 \cdot \frac{3}{8}
                                           + 2 \cdot \frac{3}{8}
                                           + 3 \cdot \frac{1}{8}
                                          = \frac{3}{8} + \frac{6}{8} + \frac{3}{8}
                                          = \frac{3}{2}.
                                    $
                              \end{ex}
                              %
                              \begin{ex}
                                    On lance un d\'e non truqu\'e.
                                    Soit $X$ la variable al\'eatoire qui est \'egale \`a $-1$
                                    si le r\'esultat est diff\'erent de $6$
                                    et $3$ si le r\'esultat est \'egal \`a $6$.
                                    On a :
                                    $E[X] = -1 \cdot \frac{5}{6} + 3 \cdot \frac{1}{6}
                                          = -\frac{5}{6} + \frac{3}{6}
                                          = -\frac{1}{3}.
                                    $
                              \end{ex}
                        %
                        \subsubsection{Variance et \'ecart-type}
\index{Variance}
\index{Ecart-type@\'Ecart-type}
                        Soit $X$ une variable al\'eatoire.
                        On d\'efinit la \emph{variance} de $X$ par
                        \par\smallskip\noindent
                        \defbox{$\displaystyle
                              \begin{aligned}
                                    Var(X)=V(X)&=E\left[(X-E[X])^2\right] \\
                                               &=E[X^2]-\left(E[X]\right)^2,
                              \end{aligned}$}
                        \par\smallskip
                        et l'\emph{\'ecart-type} par
                        \par\smallskip\noindent
                        \defbox{$\sigma(X)=\sqrt{V(X)}$}.
                        \par\smallskip
                        
                        Si la variable al\'eatoire $X$ prend les valeurs 
                        $x_1,\ x_2,\ \ldots,\ x_n$ avec les probabilit\'es 
                        $p_1,\ p_2,\ \ldots,\ p_n$, en notant $\bar{x} = E[X]$,
                        on a :
                        \par\smallskip\noindent
                        \propbox{$\displaystyle
                              \begin{aligned}
                              V(X)={}&(x_1-\bar{x})^2p_1+(x_2-\bar{x})^2p_2 \\
                                    &+\cdots+(x_n-\bar{x})^2p_n.
                              \end{aligned}$}
                        \par\smallskip
                        %
                  \subsection{Loi binomiale $B(n;p)$}
\index{Binomiale!Loi}
                  On consid\`ere $n$ \'epreuves de Bernoulli avec une probabilit\'e $p$ de succ\`es.
                  Soit $X$ la variable al\'eatoire correspondant au nombre de succ\`es.
                  L'univers image de $X$ est alors 
                  $\{ 0,\ 1,\ \ldots,\ n\}$.
                  On sait que $P(X = k) = \C{n}{k}\ p^k \cdot (1-p)^{n-k}$.
                  
                  On a 
                  \begin{align*}
                        E[X] & = \sum_{k=0}^{n} k \cdot P(X=k)                                                
                               = \sum_{k=0}^{n} k \C{n}{k} \cdot p^k \cdot (1-p)^{n-k}                        \\
                             & = \sum_{k=1}^{n} k \cdot \frac{n!}{k! (n-k)!} p^k (1-p)^{n-k}                  \\
                             & = \sum_{k=1}^{n} \frac{n \cdot (n-1)!}{(k-1)! \left[ (n-1) - (k-1) \right]!} 
                                       p \cdot p^{k-1} \cdot (1-p)^{(n-1) - (k-1)}                            \\
                             & = n p \sum_{k=1}^{n} \frac{(n-1)!}{(k-1)! \left[ (n-1) - (k-1) \right]!}
                                               p^{k-1} \cdot (1-p)^{(n-1) - (k-1)}                            \\
                             & = n p \sum_{\ell=0}^{m} \frac{m!}{\ell! (m-\ell)!} p^{\ell} \cdot (1-p)^{m-\ell}               
                               = n p \left(p+(1-p)\right)^{m}                                                 
                               = n p \cdot 1^{m}                                                              \;.
                  \end{align*}
                  \thmbox{$E[X] = n p $}.
                  
                  On peut aussi d\'emontrer que 
                  \thmbox{$V(X) = n p (1 - p) $} et donc \thmbox{$ \sigma (X) = \sqrt{n p (1 - p)}$}.
                        














\TCDLessonStyleOff
\begin{fiche}
	\begin{prop}
		\begin{itemize}
			\item Dans le cas o\`u on a \'equiprobabilit\'e, la probabilit\'e d'un \'ev\`enement est
				\thmbox{$\frac{\text{nombre de cas possibles}}{\text{nombre de cas favorables}}$}.
%				\begin{ex}
%					On consid\`ere un d\'e icosa\`edrique (\`a 20 faces) non-truqu\'e dont les faces sont num\'erot\'ees de 1 \`a 20.
%					On lance une fois ce d\'e, on consid\`ere les \'ev\`enements :
%					\begin{description}
%						\item $A$ : << Avoir un nombre premier >>,
%						\item $B$ : << Avoir un nombre inf\'erieur ou \'egal \`a 15 >>.
%					\end{description}
%					%
%					\begin{valiny}
%						\begin{description}
%							\item Les nombres premiers inf\'erieurs \`a 20 sont : 2, 3, 5, 7, 11, 13, 17 et 19 ; il y en a 8.
%								Ainsi, $P(A) = \frac{8}{20} = \frac{2}{5}$.
%							\item $P(B) = \frac{15}{20} = \frac{3}{4}$.
%						\end{description}
%					\end{valiny}
%				\end{ex}
			\item 
				La probabilit\'e d'un \'ev\`enement est \'egal \`a la somme des \'ev\`enements \'el\'ementaires qui le constituent :
				$\displaystyle{P(A) = \sum_{a \in A} P({a})}$.
%					\begin{ex}
%						On consid\`ere un d\'e cubique t\'etra\`edrique (\`a 4 faces) truqu\'e de mani\`ere \`a ce que : 
%							$p_1 = \frac{1}{2}$, $p_2 = \frac{1}{3}$, $p_3 = \frac{1}{6}$ et $p_4 = 0$ ;
%							$p_i$ d\'esigne la probabilit\'e que la face num\'erto\'ee $i$ soit cach\'ee.
%						Soient les \'ev\`enements : 
%						\begin{description}
%							\item $A$ : << Le nombre cach\'e est pair >>.
%							\item $B$ : << Le nombre cach\'e est premier >>.
%						\end{description}
%						%
%						\begin{valiny}
%							Alors 
%							\begin{description}
%								\item $P(A) = p_2 + p_4 = \frac{1}{3} + 0 = \frac{1}{3}$,
%								\item $P(B) = p_2 + p_3 = \frac{1}{3} + \frac{1}{6} = \frac{1}{2}$.
%							\end{description}
%						\end{valiny}
%					\end{ex}
			\item 
				$\displaystyle{P(\overline{A})} = 1 - P(A)$,
%				\begin{ex}
%					On lance un d\'e cubique dont les faces sont num\'erot\'es de 1 \`a 6.
%					Le d\'e est truq\'e de mani\`ere \`a ce que : 
%					$p_1 = \frac{1}{4}$, 
%					$p_2 = \frac{1}{4}$,
%					$p_3 = \frac{1}{8}$,
%					$p_4 = \frac{1}{8}$,
%					$p_5 = \frac{1}{8}$,
%					$p_6 = \frac{1}{8}$,
%					Soient les \'ev\`enements : 
%					\begin{description}
%						\item $A$ : << Avoir un num\'ero sup\'erieur ou \'egal \`a 2 >>,
%						\item $B$ : << Avoir un num\'ero inf\'erieur \`a 6 >>.
%					\end{description}
%					%
%					\begin{valiny}
%					Consid\'erons les \'ev\`enements contraires : 
%						\begin{description}
%							\item $\overline{A}$ : << Obtenir le num\'ero 1 >>,
%							\item $\overline{B}$ : << Obtenir le num\'ero 6 >>.
%						\end{description}
%					Comme 
%						\begin{description}
%							\item $P(\overline{A}) = \frac{1}{4}$, alors $P(A) = 1 - \frac{1}{4} = \frac{3}{4}$,
%							\item $P(\overline{B}) = \frac{1}{8}$, alors $P(B) = 1 - \frac{1}{8} = \frac{7}{8}$.
%						\end{description}
%					\end{valiny}	
%				\end{ex}
			\item Si $A$ et $B$ sont ind\'ependants, alors 
				\thmbox{$P(A \cap B) = P(A) \times P(B)$}.
%				\begin{ex}
%					On dispose d'un d\'e octa\`edrique (\`a 8 faces) normal dont les faces sont num\'erot\'es de 1 \`a 8.
%					On lance ce d\'e deux fois de suite et de mani\`ere ind\'ependantes.
%					
%					D\'eterminer la probabilit\'e des \'ev\`enements : 
%					\begin{description}
%						\item $A$ : << Avoir 1 au premier tirage et un nombre impair au second tirage >>,
%						\item $B$ : << Avoir un nombres premier au second tirage >>.
%					\end{description}
%					%
%					\begin{valiny}
%						\begin{description}
%							\item Pour l'\'ev\`enement $A$ :
%								La probabilit\'e d'avoir 1 au premier tirage est $\frac{1}{8}$ tandis que 
%								la probabilit\'e d'avoir un nombre impair au second tirage est $\frac{4}{8} = \frac{1}{2}$.
%								Les deux tirages \'etant ind\'ependants :
%								$P(A) = \frac{1}{8} \times \frac{1}{2} = \frac{1}{16}$.
%							\item Pour l'\'ev\`enement $B$ :
%								Il suffit juste d'avoir un nombre au premier tirage, c'est \`a dire avec probabilit\'e $1$.
%								La probabilit\'e d'avoir un nombre premier au second tirage est $\frac{4}{8} = \frac{1}{2}$.
%								$P(B) = 1 \times \frac{1}{2} = \frac{1}{2}$.
%						\end{description}
%					\end{valiny}
%				\end{ex}
		\end{itemize}
	\end{prop}
\end{fiche}



























\TCDLessonStyleOff
\clearpage
\section{Exercices}

\providecommand{\EUR}[1]{#1~\texteuro}

\subsection*{Vocabulaire}
\begin{exo}
Dans chacune des situations décrites ci-dessous, énoncer l’événement contraire de l’événement donné.
\begin{enumerate}
\item Dans une classe, on choisit deux élèves au hasard. $A$ : <<~Les deux élèves sont des filles~>>.
\item Dans un groupe de Chinois et d'Indiens, on discute avec une personne. $B$ : <<~La personne est une femme chinoise~>>.
\item Au restaurant, Luc prend un plat et un dessert. $C$ : <<~Luc prend une viande et une glace~>>.
\item À une loterie, Élise achète 3 billets.
      $D$ : <<~L’un des billets au moins est gagnant~>> , $E$ : <<~Deux billets au maximum sont gagnants~>>.
\end{enumerate}
\end{exo}
%
\begin{exo}
Une urne contient des boules blanches, noires et rouges. On tire une boule de l’urne. On note : \\
$A$ : <<~Tirer une boule blanche~>>.                  \\
$B$ : <<~Tirer une boule ni blanche ni rouge~>>.      \\
$C$ : <<~Tirer une boule noire ou une boule rouge~>>. 
\begin{enumerate}
\item $A$ et $B$ sont-ils incompatibles\,?
\item $B$ et $C$ sont-ils incompatibles\,?
\item Traduire par une phrase ne comportant pas de négation $\bar{A}$ et $\bar{B}$ .
\end{enumerate}
\end{exo}
%
\begin{exo}
Lors d’un jet de deux dés cubiques, on s’intéresse aux événements suivants : \\
$A$ : <<~La somme obtenue est au moins égale à 5~>>.                  \\
$B$ : <<~La somme obtenue est au plus égale à 5~>>.                   \\
$C$ : <<~La somme obtenue est strictement inférieure à 3~>>. 
\begin{enumerate}
\item $A$ et $B$ sont-ils contraires\,?
\item $\bar{B}$ et $C$ sont-ils incompatibles\,?
\item Traduire par une phrase $\bar{C}$.
\item $A$ et $\bar{C}$ sont-ils incompatibles\,?
\end{enumerate}
\end{exo}
%
\subsection*{D\'enombrement et probabilit\'e}
\begin{exo}
On choisit une carte au hasard dans un jeu de 32 cartes. On note : \newline
$A$ l'événement : <<~La carte choisie est un pique~>>.                     \\
$B$ l'événement : <<~La carte choisie est rouge (cœur ou carreau)~>>.      \\
$C$ l'événement : <<~La carte choisie est une figure (valet, dame, roi)~>>.
\begin{enumerate}
\item Présenter un modèle mathématique décrivant l’expérience aléatoire 
      (D\'eterminer $\Omega$ et $\card{\Omega}$ ; y a-t-il \'equiprobabilit\'e\,?). \\
\item Déterminer les probabilités des évènements $A$, $B$, $C$, $A\cap B$, $B\cap C$, $A\cup B$, $A\cup C$.
\item Déterminer la probabilité de l'événement $D$ <<~La carte choisie n'est ni un pique ni une figure~>>.
\end{enumerate}
\end{exo}
%
\begin{exo}
On jette une pièce de monnaie 3 fois de suite.
\begin{enumerate}
\item Donner la liste de tous les résultats possibles en notant $P$ pour Pile et $F$ pour Face (exemple : $PPF$).
\item Donner la probabilité des événements suivants : \newline
      $A$ <<~le tirage ne comporte que des Piles~>>.       \\
      $B$ <<~le tirage comporte au moins une fois Face~>>. 
\end{enumerate}
\end{exo}
%
\begin{exo}
Dans une assemblée de 250 personnes, on remarque que certains hommes portent une cravate et certains hommes ont les yeux marrons. Il y a 120 hommes qui portent une cravate, 85 hommes qui ont les yeux marrons, dont 50 portent une cravate.
On discute avec une personne choisie au hasard dans cette assemblée.
\begin{enumerate}
\item Quelle est la probabilité que ce soit un homme portant une cravate\,?
\item Quelle est la probabilité que ce soit un homme aux yeux marrons et portant une cravate\,?
\item Quelle est la probabilité que ce soit un homme aux yeux marrons ou portant une cravate\,?
\item Quelle est la probabilité de discuter avec une personne qui n’est ni un homme aux yeux marrons, ni un homme portant une cravate\,?
\end{enumerate}
\end{exo}
%
\begin{exo}
Lors d’un référendum, deux questions étaient posées.
65 \% des personnes ont répondu <<~oui~>> à la première question, 51 \% ont répondu <<~oui~>> à la seconde question, et 46 \%
ont répondu <<~oui~>> aux deux questions.
\begin{enumerate}
\item Quelle est la probabilité qu’une personne ait répondu <<~oui~>> à l’une ou l’autre des questions\,?
\item Quelle est la probabilité qu’une personne ait répondu <<~non~>> aux deux questions\,?
\end{enumerate}
\end{exo}
%
\begin{exo}
On lance un dé à 6 faces. On note $p_i$ la probabilité de sortie de la face marquée $i$. 
Ce dé est truqué de telle sorte que les probabilités de sortie des faces sont :
       $p_1 = 0,1$ ; $p_2 = 0, 2$ ; $p_3 = 0,3$ ; $p_4 = 0,1$ ; $p_5 = 0,15$. \newline
Quelle est la probabilité de sortie de la face marquée 6\,? \\
Quelle est la probabilité d’obtenir un nombre pair\,?
\end{exo}
%
\begin{exo}
On lance un dé à 6 faces. 
On suppose que la probabilité d’apparition de chaque face est proportionnelle au numéro inscrit sur elle 
(c'est-\`a-dire que la probabilit\'e d'obtenir la face $i$ est \'egale \`a $k \cdot i$ pour un certain $k$ ne d\'ependant pas de $i$).
Calculer la probabilité d’apparition de chaque face.
Calculer la probabilité d’obtenir un nombre pair.
\end{exo}
%
\subsection*{Probabilité conditionnelles.}
%
\subsection*{Loi binomiale}
%
\subsection*{Variable al\'eatoire}

